% Additional explanations for the final group of lessons.

\expandafter\def\csname sourcewalkthrough@50\endcsname{%
The restriction against division matters most when zeros appear. Prefix and
suffix products avoid that problem completely. First move left to right. At
index \(i\), store the product of every value strictly before \(i\), then
multiply the running prefix by \texttt{nums[i]}.

Next move right to left with a running suffix. Multiply the stored prefix by
the product of every value strictly after the index, then update the suffix.
For \texttt{[1,2,3,4]}, the first pass stores \texttt{[1,1,2,6]}. The second
pass multiplies by suffixes 24, 12, 4, and 1 to produce
\texttt{[24,12,8,6]}. One zero makes every answer except its own position
zero; two zeros make every answer zero, all without a special branch. The
returned array is not counted as auxiliary space, so the algorithm uses
constant extra state.}

\expandafter\def\csname sourcewalkthrough@51\endcsname{%
The main difficulty is finding the node to remove without first measuring the
list length. Put both pointers at a dummy node before the head. Move the fast
pointer \(n+1\) steps so there are exactly \(n\) nodes between fast and slow.

Then move both pointers together until fast becomes null. Slow now points to
the node immediately before the target, so one link update removes it. The
dummy node is important when the first real node must be deleted. For
\texttt{1->2->3->4->5} with \(n=2\), slow stops at 3 and changes its next
link from 4 to 5. The method assumes \(n\) is valid for the list. It uses one
pass after creating the gap, runs in \(O(L)\) time, and uses constant extra
space.}

\expandafter\def\csname sourcewalkthrough@52\endcsname{%
Reordering is easier when split into three familiar linked-list operations.
First find the middle with slow and fast pointers. Cut the list into two
halves. Reverse the second half, then weave one node from the first half with
one node from the reversed half.

For \texttt{1->2->3->4->5}, the halves become \texttt{1->2->3} and
\texttt{4->5}. Reversing the second gives \texttt{5->4}. Weaving produces
\texttt{1->5->2->4->3}. Before changing either next pointer, save both next
nodes so the unprocessed suffixes are not lost. Odd and even lengths differ
only in whether the first half has one extra node. No values need to be copied
and no array is required. Each phase is linear, so the complete operation is
\(O(n)\) time and \(O(1)\) extra space.}

\expandafter\def\csname sourcewalkthrough@53\endcsname{%
Reversing a list means redirecting each next pointer toward the node that was
previously before it. Three pointers are enough: \texttt{previous},
\texttt{current}, and a saved \texttt{next} node.

At each step, save \texttt{current->next} before overwriting it. Point current
backward to previous, then advance previous and current. Starting with
\texttt{1->2->3}, the first update makes \texttt{1->null}; the second makes
\texttt{2->1}; the third makes \texttt{3->2->1}. When current becomes null,
previous is the new head. Forgetting to save the forward link loses the rest
of the list after the first update. The empty list and a one-node list pass
through the same loop without special restructuring. Every node is changed
once, giving linear time and constant extra space.}

\expandafter\def\csname sourcewalkthrough@54\endcsname{%
The result has exactly 32 bits, so the loop should always run 32 times,
including leading zeros. On each iteration, shift the result left to make room
for one new bit. Copy the lowest bit of the input with \texttt{n \& 1}, then
shift the input right.

This reads the input from right to left while building the answer from left to
right. For a small pattern such as \texttt{00000101}, the read bits are
1, 0, 1, followed by zeros, and they appear at the opposite end of the result.
Use an unsigned 32-bit type so the right shift inserts zeros instead of
extending a sign bit. A loop that stops when \(n\) becomes zero would fail to
place the remaining leading zeros in their correct reversed positions unless
the result position were tracked separately.}

\expandafter\def\csname sourcewalkthrough@55\endcsname{%
A square matrix can be rotated without another matrix by separating the
movement into two simple operations. Transpose across the main diagonal, then
reverse every row. The transpose changes rows into columns; the row reversal
puts those columns in clockwise order.

For
\(\begin{smallmatrix}1&2&3\\4&5&6\\7&8&9\end{smallmatrix}\), transposing gives
\(\begin{smallmatrix}1&4&7\\2&5&8\\3&6&9\end{smallmatrix}\). Reversing each
row gives
\(\begin{smallmatrix}7&4&1\\8&5&2\\9&6&3\end{smallmatrix}\). During the
transpose, swap only cells above the diagonal with their reflected partners;
visiting both halves would swap every pair twice and undo the work. The matrix
must be square for this direct in-place mapping. Time is \(O(n^2)\), with
constant auxiliary space.}

\expandafter\def\csname sourcewalkthrough@56\endcsname{%
Two trees are the same only when their structure and values match at every
corresponding position. The recursive cases follow that statement directly.
Two null pointers match. One null pointer and one real node do not. Two real
nodes match only if their values match and both pairs of children match.

The recursion should compare left with left and right with right. It cannot
ignore missing children, because trees with the same values in different
shapes are not the same. The method stops as soon as a mismatch is found.
When the trees match, every corresponding node is visited once, so time is
\(O(n)\). The call stack uses \(O(h)\) space, where \(h\) is the larger tree
height. Testing two empty trees, one empty tree, unequal roots, and equal
values with different shapes covers the important base cases.}

\expandafter\def\csname sourcewalkthrough@57\endcsname{%
Serialization needs to preserve structure, not only values. Preorder traversal
records a node before its children, and a special null marker records every
missing child. Those markers make the representation unambiguous.

For a node, write its value, then serialize the left and right subtrees. For a
null pointer, write a marker such as \texttt{N}. Deserialization reads tokens
in the same order. A null marker returns a null pointer; a number creates a
node whose left and right children are built recursively. A tree containing
only values cannot distinguish a left child from a right child, which is why
the markers are essential. Negative and multi-digit values need a delimiter,
not character-by-character parsing. Both directions process one token per
node or null link, so their time is linear in the representation size.}

\expandafter\def\csname sourcewalkthrough@58\endcsname{%
Using the first row and first column as marker storage avoids a second matrix.
When a zero is found at \((r,c)\), set the first cell of row \(r\) and the
first cell of column \(c\) to zero. Later passes use those markers to clear
interior cells.

The first row and first column need separate Boolean flags because their cells
are also being used as storage. Scan them before writing markers. Then mark
zeros from the interior, clear interior rows and columns from the markers, and
finally clear the first row or column if its original flag was true. Changing
cells to zero during the first scan without this plan can create new zeros
that incorrectly spread farther. Every cell is visited a constant number of
times, giving \(O(mn)\) time and constant auxiliary space.}

\expandafter\def\csname sourcewalkthrough@59\endcsname{%
Four boundaries describe the unvisited rectangle: top, bottom, left, and
right. Read the top row from left to right and move top down. Read the right
column from top to bottom and move right left. Then, if rows remain, read the
bottom row backward; if columns remain, read the left column upward.

The two safety checks are necessary for a single remaining row or column.
Without them, those cells are added twice. For a \(3\times4\) matrix, the
first round consumes the complete outer ring, and the boundaries then describe
only the inner row. Each cell enters the output exactly once, so time is
\(O(mn)\). Apart from the returned vector, only four indexes are stored.}

\expandafter\def\csname sourcewalkthrough@60\endcsname{%
At each node of the main tree, ask whether a tree starting there is exactly
the target subtree. This uses two recursive functions with different jobs.
The outer function searches possible starting nodes; the inner function checks
whether two trees match completely.

If the inner comparison succeeds, return true. Otherwise search the current
node's left and right subtrees. A null target is a subtree of any tree, while
a non-null target cannot be found inside a null main tree. The equality helper
must check both values and structure, just like Same Tree. In the worst case,
many candidate roots have the target's root value and the comparison repeats,
leading to \(O(mn)\) time. The direct method is still clear and reliable for
interview constraints; subtree hashing is a more advanced optimization.}

\expandafter\def\csname sourcewalkthrough@61\endcsname{%
Binary addition can be expressed without the plus operator. XOR adds two bits
without carrying: equal bits produce 0 and different bits produce 1. AND finds
positions where both bits are 1, and shifting that result left places the
carry in the next position.

Repeat with \texttt{sum=a\char94 b} and \texttt{carry=(a\&b)<<1}. Replace \texttt{a}
with sum and \texttt{b} with carry until no carry remains. For 5 and 3,
\texttt{0101 XOR 0011} gives \texttt{0110}, while the carry is
\texttt{0010}. Further rounds propagate that carry to produce 8. Unsigned
intermediate values avoid undefined behavior when a carry enters the sign bit;
the final bit pattern can then be converted back to \texttt{int}. Fixed-width
carry propagation takes a bounded number of rounds.}

\expandafter\def\csname sourcewalkthrough@62\endcsname{%
First count how many times each value appears. The answer needs only the
\(k\) largest frequencies, so a min heap of size at most \(k\) is enough.
Push each value-frequency pair. Whenever the heap grows beyond \(k\), remove
the smallest frequency.

After all pairs are processed, every item left in the heap belongs to the top
\(k\). The returned order usually does not matter. If there are \(u\) unique
values, counting takes expected \(O(n)\) time and heap work costs
\(O(u\log k)\). A heap containing all unique values would also work but would
use more space and perform more expensive operations. Bucket sorting by
frequency is a linear alternative, but the bounded heap is a useful pattern
whenever only a small number of best items are requested.}

\expandafter\def\csname sourcewalkthrough@63\endcsname{%
Each grid position can be reached only from the cell above or the cell to the
left. Therefore the number of paths to a cell is the sum of those two counts.
The first row and first column each contain only one path because movement is
restricted to right and down.

A one-dimensional array can represent the current row. Initialize every entry
to one. For each later row, update from left to right with
\texttt{dp[col] += dp[col-1]}. Before the update, \texttt{dp[col]} is the
count from above; \texttt{dp[col-1]} is the newly computed count from the
left. For a \(3\times7\) grid the final value is 28. This reduces space from
\(O(mn)\) to \(O(n)\) while keeping the same \(O(mn)\) time.}

\expandafter\def\csname sourcewalkthrough@64\endcsname{%
Two strings are anagrams when every character appears the same number of
times. With lowercase English letters, a fixed array of 26 counts is simpler
than sorting. First reject different lengths, because no count adjustment can
repair that mismatch.

For each position, increase the count for the character in the first string
and decrease the count for the character in the second. At the end, every
entry must be zero. Combining the two updates in one loop is safe because the
strings have equal length. For \texttt{anagram} and \texttt{nagaram}, the same
letters cancel even though their positions differ. Sorting would take
\(O(n\log n)\); counting takes \(O(n)\) time and constant alphabet space. A
larger or Unicode alphabet would use a hash map instead of a 26-cell array.}

\expandafter\def\csname sourcewalkthrough@65\endcsname{%
Move one pointer from each end. Skip any character that is not a letter or
digit, then compare lowercase forms of the remaining characters. A mismatch
proves the string is not a palindrome; matching characters let both pointers
move inward.

For \texttt{A man, a plan, a canal: Panama}, spaces and punctuation are
ignored, leaving matching pairs from the outside inward. Character functions
should receive an \texttt{unsigned char} value so non-ASCII negative
\texttt{char} values do not cause undefined behavior. The pointers must be
checked after each skip so an input containing only punctuation does not read
outside the string. Every position is passed at most once, giving linear time
and constant extra space without constructing a filtered copy.}

\expandafter\def\csname sourcewalkthrough@66\endcsname{%
An opening delimiter creates work that must be completed later, so store it on
a stack. A closing delimiter must match the most recent unmatched opening
delimiter. This last-in, first-out rule is exactly the stack order.

Push \texttt{(}, \texttt{[}, or \texttt{\{} when it appears. For a closing
character, first ensure the stack is not empty, then compare it with the top.
Pop only after a correct match. At the end, the stack must be empty; otherwise
some opening delimiter was never closed. The string \texttt{([\{\}])} works
because closures occur in reverse opening order, while \texttt{([)]} fails at
the first closing parenthesis. Each character causes at most one push or pop,
so the solution is linear.}

\expandafter\def\csname sourcewalkthrough@67\endcsname{%
Every node in a BST must satisfy restrictions created by all its ancestors,
not only by its parent. Pass a lower and upper bound into the recursive call.
The current value must lie strictly between them.

The left child inherits the lower bound and receives the current value as its
new upper bound. The right child inherits the upper bound and receives the
current value as its new lower bound. Wide integer bounds avoid trouble when a
node contains \texttt{INT\_MIN} or \texttt{INT\_MAX}. Strict comparisons also
reject duplicates when the problem defines a strict BST. A local parent check
can miss a value deep in the left subtree that is larger than the root; the
carried interval correctly rejects it. Every node is checked once.}

\expandafter\def\csname sourcewalkthrough@68\endcsname{%
Let \texttt{dp[i]} mean that the prefix ending just before index \(i\) can be
split into dictionary words. Start with \texttt{dp[0]=true}, representing the
empty prefix. From each reachable position, test dictionary words or later end
positions and mark a new prefix when the intervening substring is a word.

For \texttt{leetcode} with \texttt{leet} and \texttt{code}, the reachable
positions become 0, 4, and 8. A word match from an unreachable position must
not be used, because the text before it has no valid split. A hash set gives
fast dictionary lookup, though substring construction still has a cost. The
basic index-based solution is often described as \(O(n^2)\) plus substring
work. The key idea is that each Boolean summarizes all possible segmentations
of one prefix.}

\expandafter\def\csname sourcewalkthrough@69\endcsname{%
Searching separately from every grid cell repeats prefix work for every word.
A Trie shares those prefixes. Insert all words, then start backtracking from
each cell whose character is a child of the Trie root.

During DFS, the current Trie node tells which next letters can still lead to a
word. A missing child stops the branch immediately. When a Trie node stores a
complete word, add it to the answer and clear or mark that stored word so the
same result is not returned twice. Temporarily mark the board cell as visited,
explore four neighbors, then restore it for other paths. Removing exhausted
Trie branches is an optional pruning improvement. The important distinction is
that board state prevents cell reuse in one path, while Trie state prevents
searching prefixes that belong to no requested word.}

\expandafter\def\csname sourcewalkthrough@70\endcsname{%
Try to match the word one character at a time. A recursive state contains the
grid position and the index of the next required character. If the index
reaches the word length, the complete word has been matched.

Reject a state when it is outside the grid, the cell is already used in this
path, or its character differs. Otherwise temporarily mark the cell, explore
the four neighbors for the next character, and restore the cell before
returning. Restoration is essential because a cell used by one attempted path
may be needed by another. The outer loops try every cell as a possible start.
Simple checks such as a word longer than the number of cells can fail early.
The worst case is exponential because many matching prefixes may branch, but
pruning happens immediately on every mismatch.}

\expandafter\def\csname sourcewalkthrough@71\endcsname{%
At least one half of the current interval is normally sorted, even though the
whole array is rotated. Compare the middle value with the boundaries to decide
which half has normal order. Then ask whether the target lies inside that
sorted value range.

If the left half is sorted and the target is between the left value and the
middle value, discard the right half; otherwise discard the left. Apply the
mirror rule when the right half is sorted. For \texttt{[4,5,6,7,0,1,2]} and
target 0, the first middle is 7. The left half is sorted but does not contain
0, so the search moves right and finds the target. Careful inclusive and
exclusive comparisons prevent losing a boundary target. Each decision halves
the search range, so time is logarithmic when values are distinct.}

\expandafter\def\csname sourcewalkthrough@72\endcsname{%
The direct dynamic-programming state is \texttt{dp[i]}, the length of the
longest increasing subsequence that ends exactly at index \(i\). Every single
value can form a length-one subsequence. Look at earlier positions \(j<i\);
when \texttt{nums[j] < nums[i]}, value \(i\) can extend the subsequence ending
at \(j\).

Take the largest \texttt{dp[j]+1} among those valid predecessors, then keep the
largest state anywhere as the final answer. In \texttt{[10,9,2,5,3,7,101,18]},
the value 7 can extend subsequences ending at 2, 5, or 3; the best choice gives
length three at that point. The values need not be adjacent, which separates a
subsequence from a subarray. This clear recurrence uses \(O(n^2)\) time and
\(O(n)\) space; the tails/binary-search method is a later optimization.}

\expandafter\def\csname sourcewalkthrough@73\endcsname{%
To reach step \(i\), the final move must come from step \(i-1\) or step
\(i-2\). Therefore the number of ways is the sum of the previous two counts.
This is the Fibonacci pattern with problem-specific starting values.

For one step there is one way. For two steps there are two ways: \(1+1\) or
2. Keep only those two previous values because no older state is needed. Each
new step computes their sum, then shifts the pair forward. For five steps the
counts are 1, 2, 3, 5, and 8. Defining the base cases clearly prevents the
common off-by-one error. The loop uses \(O(n)\) time and \(O(1)\) space,
whereas plain recursion repeats the same smaller stair counts many times.}

\expandafter\def\csname sourcewalkthrough@74\endcsname{%
Inverting a tree is a local operation repeated at every node. Swap the left
and right child pointers, then recursively invert the two child subtrees. A
null pointer is the base case.

The order can be swap-then-recurse or recurse-then-swap as long as the calls
follow the pointers consistently. For a root 4 with children 2 and 7, the
first swap places 7 on the left and 2 on the right; recursion then mirrors the
children below them. No new nodes are required because only links change.
Every node is visited once, so time is \(O(n)\). Recursive stack space is
\(O(h)\); a queue can perform the same transformation breadth first. A
one-node tree and a completely skewed tree are useful tests.}

\expandafter\def\csname sourcewalkthrough@75\endcsname{%
Water normally flows from a cell to an equal or lower neighbor. Starting from
every cell and simulating that direction repeats a great deal of work. Reverse
the question: begin at each ocean and move from a cell to neighbors with equal
or greater height. Every reached cell can flow back down to that ocean.

Run one multi-source DFS or BFS from the Pacific borders and another from the
Atlantic borders. Store the reached cells in two Boolean grids. A cell belongs
in the answer when both grids mark it. The reverse height comparison is the
most important detail; using the original downhill comparison from the ocean
would move in the wrong direction. Each ocean traversal processes a cell at
most once, so total time and storage are \(O(mn)\). Corners may begin in both
searches, and equal-height plateaus are connected in both directions.}